% Solution: Mas-Colell, Whinston & Green (1995) Exercise 9.B.4. Input by \solitem.
% New in the 50-item list. Check: check_B_L03.py. Compared with the MWG solutions manual.
\begin{exsol}{MWG Exercise 9.B.4}{302}{Advanced}{4}
\asked Prove that in a finite zero-sum game of perfect information all subgame-perfect Nash
equilibria give the same payoffs.

\method Backward induction and Zermelo's theorem (MWG \S9.B, Prop.~9.B.1, p.~272; Lecture~3,
slide~30, whose corollary gives uniqueness only when no player has equal payoffs at two terminal
nodes); subgames and SPNE (MWG Defs.~9.B.1--9.B.2, pp.~274--275; Lecture~3, slides~32--34); in a
finite game of perfect information the SPNE are exactly the backward-induction profiles (MWG
Prop.~9.B.2; Lecture~3, slide~35). The proof is an induction on subgames that tracks
\emph{values}, not strategies: ties may create many SPNE, and the point is that they cannot create
different payoffs.

\solution
\textbf{Setting.} Two players; a finite tree in which every information set is a single node;
$u_1(z) + u_2(z) = 0$ at every terminal node $z$, so we can track $v(z) = u_1(z)$ and $u_2(z) =
-v(z)$. For a decision node $x$, $\Gamma_x$ is the subgame starting at $x$ (every node starts one,
by perfect information), and $C(x)$ is the set of its immediate successors. An SPNE may use
behavioral strategies.

\textbf{Claim.} For every node $x$ there is a number $V(x)$ such that \emph{every} SPNE of
$\Gamma_x$ gives player~1 the payoff $V(x)$ (and player~2 $-V(x)$).

\textbf{Proof by induction on the length of the longest path from $x$ to a terminal node.}

\emph{Base.} If $x$ is terminal, $V(x) = v(x)$.

\emph{Step.} Let $x$ be a decision node of player~$i$ and suppose the claim holds for every
$y \in C(x)$. Take any SPNE $\sigma$ of $\Gamma_x$.
\begin{enumerate}[label=(\roman*)]
  \item For each $y \in C(x)$, the restriction of $\sigma$ to $\Gamma_y$ is an SPNE of $\Gamma_y$ (a
  subgame of $\Gamma_y$ is a subgame of $\Gamma_x$). By the induction hypothesis it gives player~1
  exactly $V(y)$. So if play moves from $x$ to $y$, the continuation pays $(V(y), -V(y))$, whatever
  $\sigma$ prescribes inside $\Gamma_y$.
  \item The restriction of $\sigma$ to $\Gamma_x$ is a Nash equilibrium of $\Gamma_x$. Player~$i$ could
  change only his action at $x$, keeping the rest of $\sigma$; moving to $y$ would give him
  $V(y)$ if $i = 1$, or $-V(y)$ if $i = 2$. Therefore every action played with positive probability at
  $x$ must lead to a maximizer:
  \[
    i = 1:\ \ y \in \arg\max_{y' \in C(x)} V(y'), \qquad
    i = 2:\ \ y \in \arg\max_{y' \in C(x)} \bigl(-V(y')\bigr) = \arg\min_{y' \in C(x)} V(y').
  \]
  \item All the successors in the support therefore have the same value, so the payoff of $\sigma$ in
  $\Gamma_x$ is
  \[
    V(x) = \max_{y \in C(x)} V(y) \ \text{ if } i = 1, \qquad V(x) = \min_{y \in C(x)} V(y) \ \text{ if } i = 2,
  \]
  a number that does not depend on which SPNE $\sigma$ we took.
\end{enumerate}
Applying the claim at the root: every SPNE of the game gives $(V(x_0), -V(x_0))$. SPNE exist by
Zermelo's theorem, so the set of SPNE payoffs is exactly $\{(V(x_0), -V(x_0))\}$. $\square$

\textbf{Where zero-sum is used.} In step~(iii). When player~$i$ is indifferent among several
successors, they give \emph{him} the same payoff, and in a two-player zero-sum game they then give
the opponent the same payoff too. Without zero-sum, the opponent's payoff can differ across the
tied successors. Counterexample: player~1 chooses $l \mapsto (0,1)$ or $r \mapsto (0,2)$; both
$l$ and $r$ are SPNE, with different payoffs for player~2.

\textbf{Two remarks.}
\begin{itemize}
  \item \emph{Two players are needed.} With three players and $u_1 + u_2 + u_3 = 0$ the result fails:
  player~1 choosing between $(0, 1, -1)$ and $(0, -1, 1)$ has two SPNE with different payoffs for
  players~2 and~3. ``Zero-sum'' in the exercise means two-player zero-sum.
  \item \emph{Chance moves.} If the tree has moves of Nature with known probabilities, define $V(x)$ at
  a chance node as the expected value of its successors' values; the same induction goes through.
  By the minimax theorem $V(x_0)$ is also the value of the game, and every Nash equilibrium, not
  only every SPNE, yields it.
\end{itemize}

\interp In a pure conflict, ties are harmless: if I am indifferent between two moves, so are you,
so my tie-breaking cannot be used to help or hurt you. That is why chess-like games have a
well-defined value even though the backward-induction strategies are far from unique. In the
non-zero-sum games of Lecture~4 (entry, Brexit) a follower's indifference \emph{is} strategically
important: how she breaks it can change the leader's payoff (compare MWG 9.B.5, item ``MWG 9.B.5'',
part (e)).

\verify \texttt{check\_B\_L03.py} draws 400 random two-player zero-sum trees (depth up to 4, two or
three moves per node, integer payoffs in $\{-2, \dots, 2\}$, fixed seed), computes the set of pure
SPNE payoff vectors over all tie-breakings, and finds exactly one in each tree; 151 of the trees have
more than one SPNE, so ties do occur. It also confirms the two counterexamples (non-zero-sum and
three-player zero-sum) have two SPNE payoff vectors. Output line: \texttt{mwg-9-b-4: 400 random
zero-sum trees, unique SPNE payoffs (151 with several SPNE)}.
\emph{Comparison with the MWG solutions manual} (Exercise 9.B.4, p.~9-2 of the manual; label read
as ``9.8.4'' in the OCR): it argues that if a player has equal payoffs at two terminal nodes, the
zero-sum property makes the other player indifferent too, so ``there may exist many different
\dots{} SPNE \dots{} but all of them will yield the same payoffs''. \textbf{Agreement} on the result;
the manual's argument speaks of terminal nodes only, and the induction above extends it to ties
between subgame values, which is where it is actually needed.

\mistakes
\begin{itemize}
  \item Proving that the SPNE is unique: it need not be (ties); only the payoffs are.
  \item Using Zermelo's theorem's uniqueness clause, which requires no ties at all, instead of an
  argument that allows ties.
  \item Forgetting that SPNE may be mixed at tied nodes; the argument must cover any probability on
  the maximizers, which (ii)--(iii) do.
\end{itemize}

\answer{By induction on subgames: if every subgame after $x$ has a unique SPNE value $V(y)$ for
player~1, then at $x$ any SPNE puts weight only on successors that maximize the mover's payoff,
$\max_y V(y)$ for player~1 or $\min_y V(y)$ for player~2, so $V(x)$ is unique. At the root all SPNE
pay $(V(x_0), -V(x_0))$. Zero-sum (with two players) is what makes the mover's ties irrelevant for
the opponent.}
\end{exsol}
